\subsection{Repeated local quantum readout with controlled disturbance}
\label{sec:source-scalar-sequential-instrument}

The same finite scalar action admits an explicit instrument for its bounded
sine detector. This gives a joint classical record algebra and a controlled
comparison of sequential measurement probabilities. The construction assumes
quantum pointer preparation, controlled field operations and Born readout;
these operations are additional inputs. It does not identify the classical
field registers with sampled quantum outcomes.

Use the Hamiltonian, original vacuum, coherent intervention and symmetric
step $S$ of the finite scalar probability construction. Write
$A_g=\Phi(g)=\sum_i m_i g_i q_i$ and $W(g)=e^{iA_g}$, with $\hbar=1$.
A fresh pointer starts in $|+X\rangle$. Apply
$U_g=\exp(-iZ\otimes A_g/2)$ and read the standard Pauli $Y$, labeling
its positive eigenvalue by zero. The two outcome operators are
\[
 K_0=\frac{e^{-iA_g/2}-i e^{iA_g/2}}2,\qquad
 K_1=\frac{e^{-iA_g/2}+i e^{iA_g/2}}2.
\]
In particular $K_0^*K_0=(I+\sin A_g)/2$ and
$\sum_sK_s^*K_s=I$. The sign convention matters: reversing the sign
of $U_g$ with the same pointer labels reverses the sine response.

\begin{theorem}[Centered sequential instrument]
\label{thm:source-scalar-sequential-instrument}
Apply this instrument after every step $S$. For a record
$h=(h_1,\ldots,h_n)$ let
\[
 L_h=K_{h_n}S\cdots K_{h_1}S,\qquad
 \mathcal J_n(\rho)=\sum_h L_h\rho L_h^*\otimes|h\rangle\langle h|.
\]
This is a completely positive trace-preserving map into a field algebra
and a classical record algebra. Its finite prefixes are consistent when
only the classical records are retained. The probability of a record is
$\operatorname{Tr}(L_h\rho L_h^*)$; these records need not be independent.

Let $p_j$ be the original unmeasured probability after $j$ steps and define
$[A_g(0),A_g(d\tau)]=i c_d$ for the split evolution. Then the
unconditional probability of outcome zero at the $j$th readout is exactly
\[
 p_j^{\rm seq}-\tfrac12
 =D_j\bigl(p_j-\tfrac12\bigr),\qquad
 D_j=\prod_{d=1}^{j-1}\cos(c_d/2).                 \tag{*}
\]
The nonselective readout preserves all canonical first moments. It increases
the expected original Hamiltonian by $G_0/8$ at each readout, where
$G_0=\langle g,g\rangle_M$. Both the split and continuous sequential
vacuum baselines remain exactly $1/2$.
\end{theorem}
\begin{proof}
The bounded outcome operators give the Kraus completeness identity above;
iterating it proves normalization of every record distribution and its
prefix marginals. Summing the two outcomes cancels the interference terms:
\[
 \mathcal M_g(\rho)=\tfrac12\left(
 W(-g/2)\rho W(g/2)+W(g/2)\rho W(-g/2)\right).
\]
In the interaction picture, conjugation of the future Weyl observable by
either half kick multiplies it by $e^{\pm i c_d/2}$. Thus each earlier
nonselective readout contributes the real factor $\cos(c_d/2)$.
Taking the imaginary part of the Weyl expectation proves (*), without
assuming that the measured state remains Gaussian. A centered vacuum has
real Weyl expectation; hence its sine marginal stays zero after these
real multipliers. The half kicks shift $P$ by
$\pm M^{1/2}g/2$, leave $Q$ unchanged, and average to zero first-moment
shift. Expanding $P^TP/2$ gives the energy increment $G_0/8$.
These energy identities apply on the finite-energy domain; the instrument
and its probability identities are bounded and apply to every density state.
\end{proof}

The detector occupies the 32-site half-volume of the original finite
action. It can be compiled using qubits at those sites. Prepare their GHZ
state along a 31-edge spanning tree of the \emph{scalar action graph}, then
apply the local operations
$\exp(-iZ_i\otimes m_i g_i q_i/2)$. Read $Y$ at the root and $X$ at
the other sites, retaining every local bit. For a fine record $b$ with
parity $s$, its field outcome operator is
$2^{-(32-1)/2}K_s$. There are $2^{31}$ fine records of each parity, so
coarse-graining to parity gives exactly the instrument above.
An independent stabilizer calculation verifies the supplied tree against
the actual detector support and action edges. The instrument lies in the
original equal-time regional field algebra; it needs no exterior collar
to measure this position-only detector.

This circuit has 32 zero preparations, one Hadamard, 31 CNOTs, 32 local
controlled field kicks, 32 pointer reads and one parity decode per readout:
129 ideal operations, or 2,709 over 21 readouts. The initial coherent field
intervention uses 16 further local Weyl factors; original-vacuum preparation
is supplied separately. These counts exclude field evolution, hardware
time, finite precision and preparation of the field vacuum. The scalar
action graph is not identified with a native seam-routing implementation.
The construction supplies the controlled field couplings and entangling
links, and treats the readout as an instantaneous ideal instrument.

\begin{proposition}[Sequential comparison over recovered time intervals]
\label{prop:source-scalar-sequential-comparison}
Let $[a_j,b_j]$ be the ordered recovered time intervals of the same
reference mean-field history. Suppose the earlier probability certificate
gives $|p_j-p(t_j)|\le\epsilon_j$ for every $t_j\in[a_j,b_j]$,
and $|p_j-1/2|\le u_j$. Let $d_j\le D_j\le1$ be a positive bound and set
\[
 L_j=\frac{G_0^2}{8}\sum_{k<j}(b_j-a_k)^2<1,
 \qquad \Gamma_j=\max(1-d_j,L_j).
\]
For every ordered choice $t_k\in[a_k,b_k]$, the same continuous
Hamiltonian with the same centered instrument at those times satisfies
\[
 |p_j^{\rm seq}-p^{\rm seq}_{\rm cont}(t_1,\ldots,t_j)|
 \le\epsilon_j+\Gamma_j u_j.
\]
If the unmeasured split response has a fixed sign and magnitude at least
$s_j$, its measured split magnitude is at least $d_js_j$.
\end{proposition}
\begin{proof}
The continuous cross-time commutator is
\[
 c(t)=\langle g,\mathsf A^{-1/2}
                 \sin(t\sqrt{\mathsf A})g\rangle_M,
 \qquad |c(t)|\le G_0|t|.
\]
The same interaction-picture argument gives a continuous attenuation
product with factors $\cos(c(t_j-t_k)/2)$. The bound
$\cos x\ge1-x^2/2$ and $L_j<1$ show that all these factors are positive
and their product lies in $[1-L_j,1]$. Therefore the two attenuation
factors differ by at most $\Gamma_j$. Subtract their products with the
unmeasured responses, multiplying the original probability error by the
continuous attenuation (at most one), to obtain the stated bound.
\end{proof}

For the complete 64-mode action, the discrete commutators are computed
exactly in $\mathbb Q(\phi)$ from
\[
 z_0=0,\quad z_1=\tau g,\quad
 z_{d+1}=(2I-\tau^2\mathsf A)z_d-z_{d-1},\qquad
 c_d=\langle g,z_d\rangle_M.
\]
Independent verification instead uses first-order kick--drift--kick
phase-space evolution and the independently assembled full stiffness.
Rational bounds on $c_d^2$ give
$d_j=\prod_{d<j}(1-\overline{c_d^2}/8)$, with outward rounding retained.
All 21 times and all 210 prior-readout pairs enter the comparison. At
$j=16$, the split sequential response magnitude exceeds $0.00950581279$, its comparison
error is below $0.00083418122$, and the continuous sequential response
has the same negative sign with magnitude above $0.00867163157$.
The same 15 time slots remain resolved. Through the final readout, the
discrete attenuation loss is below $0.000044641$.

The classical clock certificate is reused as a conditional reference-time
enclosure. Centering preserves ensemble first moments; it does not make
each measurement branch obey the original deterministic configuration
history or derive a clock from the binary records. The certificate reports
unconditional marginals of an explicit joint instrument, not sampled
histories or independent trials. Original-vacuum preparation, the quantum
control and Born readout laws, the physical meaning of the clock, and
finite-duration hardware implementation remain supplied. No spatial
continuum, interacting quantum limit or observed physical outcome follows.
